\crefname{section}{Appendix}{Appendices}
\Crefname{section}{Appendix}{Appendices}

\section{Additional Related Work}
\label{app:related-work}
In this section, we discuss related work covering generative modeling, Schrödinger bridges, flow matching, and other approaches to trajectory inference.

\textbf{Generative modeling.} The machine learning community has made substantial progress in sampling from complex, unknown distributions by transporting particles using diffusion models \citep{ho2020denoising, song2019generative, song2020score}, Schrödinger bridges \citep{de2021diffusion, pavon2021data, Vargas2021, Wang2023}, continuous normalizing flows \citep{chen2018neural, grathwohl2018ffjord}, and flow matching \citep{lipman2022flow}. These methods are primarily designed for generative modeling, where the goal is to transform a simple distribution (e.g., a Gaussian) into a data distribution. While they involve a notion of “time,” it typically serves as an auxiliary dimension rather than representing physical or real-world temporal dynamics. As a result, these methods focus on two-marginal transport problems and are not naturally suited to forecasting or modeling systems with multiple observed time points.

\textbf{Trajectory inference with Schrödinger bridges, flow matching and normalizing flows.} 
Several recent works have explored the use of Schrödinger bridges (SBs) and optimal transport for modeling trajectories across time \citep{schiebinger2019optimal, yang2018scalable}. 
Most SB-based methods are restricted to pairwise interpolation between consecutive marginals and cannot handle multiple marginals in a principled way. Several extensions have been proposed to address this. \texttt{DMSB} \citep{chen2024deep} incorporates momentum into the particles to exploit local information and model multi-marginal dynamics, while \citet{hong2025trajectory} proposed generalizations using higher-order derivatives, although these approaches are computationally feasible only in low-dimensional settings. \citet{shen2024multi} introduced \texttt{SBIRR}, a multi-marginal Schrödinger bridge method with iterative reference refinement. By integrating information from all marginals into the reference dynamics, \texttt{SBIRR} enables not only interpolation but also extrapolation via learned dynamics, and serves as one of the main baselines in our experiments. \citet{tong2024simulation} introduced \texttt{SF2M}, a simulation-free SB method based on score and flow matching, designed for efficiency but limited to Brownian motion as the reference and two-marginal settings. 

\citet{hashimoto2016learning} proposed a regularized recurrent neural network trained with a Wasserstein gradient flow loss for trajectory reconstruction. \citet{bunne2022proximal} extended this direction with \texttt{JKOnet}, a neural implementation of the Jordan–Kinderlehrer–Otto (JKO) scheme, enabling learning of the underlying energy landscape. More recently, \citet{terpin2024learning} proposed \texttt{JKOnet*}, which bypasses the bilevel optimization of \texttt{JKOnet} by exploiting first-order optimality conditions of the JKO scheme, thereby enabling efficient and accurate recovery of potential, interaction, and internal energy components from population snapshots. While \texttt{JKOnet*} extends \texttt{JKOnet} by learning potential, interaction, and internal energies from population snapshots, its core assumption is that dynamics arise from gradient flow on an energy landscape. This makes it well suited for physical systems but less natural for regulatory feedback networks like the Repressilator, where oscillations are driven by transcriptional repression and delayed activation rather than energy minimization. Moreover, \texttt{JKOnet*} evaluates performance by splitting train/test data within each timepoint, whereas our setting holds out entire timepoints to assess extrapolative forecasting ability, making direct comparison challenging. 

Other works use deterministic flows to infer trajectories. \texttt{TrajectoryNet} \citep{tong2020trajectorynet} combines dynamic optimal transport with continuous normalizing flows (CNFs) to generate continuous-time, nonlinear trajectories from snapshot data. These flows are governed by ODEs rather than SDEs, and incorporate regularization that encourages short, energy-efficient paths. \citet{tong2023improving} extended flow matching by modeling marginals as mixtures indexed by a latent variable. They proposed \texttt{OT-CFM} and \texttt{SB-CFM}, where each component evolves under a simple vector field and the overall dynamics are determined by mixing strategies based on optimal transport or SB principles. Like many flow-matching methods, these approaches are restricted to two marginals and are applied piecewise for longer time series. To address issues of geometry in the data space, \citet{huguet2022manifold} proposed \texttt{MIOFlow}, which first learns the data manifold and then solves the flow problem within that learned structure. They use neural ODEs to transport points along this manifold. \citet{atanackovic2024meta} introduced methods to learn multiple flows over a Wasserstein manifold, though these too are aimed at interpolation rather than forecasting. 

Overall, these methods move particles using learned vector fields and various forms of regularization, but they are generally not designed for forecasting. Like many SB-based methods, they suffer when applied to extrapolation tasks, particularly in the absence of a good prior or reference dynamic. 

